\subsection{Rank and Invertibility Consequences}
Throughout this subsection, all matrices are over the same fixed field $\mathbb{F}$, and every matrix size is a positive integer.

\begin{corollary} \label{cor:basis-columns-form-invertible-matrix} \hlt{Basis Columns Form an Invertible Matrix}\\
Let $n\geq1$. If $(\mathbf{a}_1,\ldots,\mathbf{a}_n)$ is a basis of $\mathbb{F}^n$, then $A=[\mathbf{a}_1\ \cdots\ \mathbf{a}_n]\in\operatorname{Mat}(n,n,\mathbb{F})$ is invertible.
\begin{proof}
Let $\epsilon=(\mathbf{e}_1,\ldots,\mathbf{e}_n)$ be the ordered list of columns of $I_n$, and set $\alpha\defeq(\mathbf{a}_1,\ldots,\mathbf{a}_n)$.\\
The matrix $I_n$ is invertible with inverse $I_n$ by Proposition \ref{prop:matrix-product-associativity-and-identities}(ii), so $\epsilon$ is a basis of $\mathbb{F}^n$ by Theorem \ref{thm:invertible-matrix-columns-basis}.\\
The relation $(\forall\,j\in\{1,\ldots,n\})\,[\mathbf{a}_j]_{\epsilon}=\mathbf{a}_j$ holds by Definition \ref{def:coordinate-vector-and-coordinate-map}, Proposition \ref{prop:matrix-vector-product-column-combination}, and $I_n\mathbf{a}_j=\mathbf{a}_j$.\\
Thus $C_{\epsilon\leftarrow\alpha}=A$ by Definition \ref{def:change-of-basis-matrix} and Definition \ref{def:matrix-representation-of-linear-map}, and $A$ is invertible by Theorem \ref{thm:change-of-basis-matrices}(i).\\
\hlt{Structural Commentary}: Standard coordinates identify the column matrix with a change-of-basis matrix.\\
\hlt{Key Insight}: A basis-to-basis coordinate rewrite is invertible.
\end{proof}
\end{corollary}

\begin{corollary} \label{cor:square-matrix-invertible-iff-full-rank} \hlt{Invertibility and Full Rank}\\
Let $n\geq1$ and $A\in\operatorname{Mat}(n,n,\mathbb{F})$. Then $A$ is invertible $\Longleftrightarrow$ $\operatorname{rank}(A)=n$.
\begin{proof}
\begin{enumerate}[label=\roman*.]
\setlength{\itemsep}{0pt}
\item[$(\Rightarrow)$] Suppose $A$ is invertible. Its columns form basis of $\mathbb{F}^n$ by Theorem \ref{thm:invertible-matrix-columns-basis}; hence $\operatorname{col}(A)=\mathbb{F}^n$, $\operatorname{rank}(A)=n$ by Definition \ref{def:column-space-matrix-rank-nullity}.
\item[$(\Leftarrow)$] Suppose $\operatorname{rank}(A)=n$. Then $\dim(\operatorname{col}(A))=n=\dim(\mathbb{F}^n)$ by Definition \ref{def:column-space-matrix-rank-nullity}, so $\operatorname{col}(A)=\mathbb{F}^n$ by Proposition \ref{prop:dimension-bounds-and-basis-criteria}(ii).\\
The $n$ columns of $A$ therefore form a basis of $\mathbb{F}^n$ by Proposition \ref{prop:full-dimension-basis-criteria}(i), and $A$ is invertible by Corollary \ref{cor:basis-columns-form-invertible-matrix}.
\end{enumerate}
\hlt{Structural Commentary}: One direction reads rank from basis columns; the converse upgrades full dimension to a basis.\\
\hlt{Key Insight}: Square-matrix invertibility is equivalent to full column rank.
\end{proof}
\end{corollary}

\begin{corollary} \label{cor:transpose-invertible-iff} \hlt{Transpose Preserves Invertibility}\\
Let $n\geq1$ and $A\in\operatorname{Mat}(n,n,\mathbb{F})$. Then $A$ is invertible $\Longleftrightarrow$ $A^t$ is invertible.
\begin{proof}
\begin{enumerate}[label=\roman*.]
\setlength{\itemsep}{0pt}
\item[$(\Rightarrow)$] Suppose $A$ is invertible. Then $\operatorname{rank}(A)=n$ by Corollary \ref{cor:square-matrix-invertible-iff-full-rank}, and $\operatorname{rank}(A^t)=\operatorname{rank}(A)=n$ by Theorem \ref{thm:row-rank-equals-column-rank}.\\
Hence $A^t$ is invertible by Corollary \ref{cor:square-matrix-invertible-iff-full-rank}.
\item[$(\Leftarrow)$] Suppose $A^t$ is invertible. The matrix $(A^t)^t$ is invertible by the forward direction with $A^t$ in place of $A$.\\
Since $(A^t)^t=A$ by Proposition \ref{prop:elementary-transpose-rules}(i), $A$ is invertible.
\end{enumerate}
\hlt{Structural Commentary}: Rank invariance proves the forward direction; double transposition reverses it.\\
\hlt{Key Insight}: Transposition preserves full rank.
\end{proof}
\end{corollary}

\begin{theorem} \label{thm:complete-full-column-rank-matrix} \hlt{Completing a Full-Column-Rank Matrix}\\
Let $1\leq n<m$ and $X\in\operatorname{Mat}(m,n,\mathbb{F})$. If $\operatorname{rank}(X)=n$, then $(\exists\,X'\in\operatorname{Mat}(m,m-n,\mathbb{F}))\,\left[\begin{smallmatrix}X&X'\end{smallmatrix}\right]\in\operatorname{Mat}(m,m,\mathbb{F})^{\times}$.
\begin{proof}
Write $X=[\mathbf{x}_1\ \cdots\ \mathbf{x}_n]$. Then $\dim(\operatorname{col}(X))=n$ by Definition \ref{def:column-space-matrix-rank-nullity}. The columns form a spanning list of length $n$ in $\operatorname{col}(X)$.\\
Hence they form a basis of $\operatorname{col}(X)$ by Proposition \ref{prop:full-dimension-basis-criteria}(i) and are linearly independent in $\mathbb{F}^m$.\\
Extend them to a basis $(\mathbf{x}_1,\ldots,\mathbf{x}_m)$ of $\mathbb{F}^m$ by Theorem \ref{thm:extending-independent-lists}, and set $X'\defeq[\mathbf{x}_{n+1}\ \cdots\ \mathbf{x}_m]\in\operatorname{Mat}(m,m-n,\mathbb{F})$.\\
The matrix $\left[\begin{smallmatrix}X&X'\end{smallmatrix}\right]$ is invertible by Corollary \ref{cor:basis-columns-form-invertible-matrix}.\\
\hlt{Structural Commentary}: Column rank identifies an independent list; basis extension supplies the missing columns.\\
\hlt{Key Insight}: Full-column-rank matrices are initial blocks of invertible matrices.
\end{proof}
\end{theorem}

\begin{theorem} \label{thm:full-rank-products-preserve-rank} \hlt{Full-Rank Products Preserve Rank}\\
Let $A\in\operatorname{Mat}(m,n,\mathbb{F})$. Then:
\begin{enumerate}[label=(\roman*)]
\setlength{\itemsep}{0pt}
\item If $X\in\operatorname{Mat}(p,m,\mathbb{F})$ and $Y\in\operatorname{Mat}(n,q,\mathbb{F})$ satisfy $\operatorname{rank}(X)=m$ and $\operatorname{rank}(Y)=n$, then $\operatorname{rank}(XAY)=\operatorname{rank}(A)$.
\item If $P\in\operatorname{Mat}(m,m,\mathbb{F})$ and $Q\in\operatorname{Mat}(n,n,\mathbb{F})$ are invertible, then $\operatorname{rank}(PAQ)=\operatorname{rank}(A)$.
\end{enumerate}
\begin{proof}
\begin{enumerate}[label=(\roman*)]
\setlength{\itemsep}{0pt}
\item The equality $\dim(\ker(X))=0$ holds by Corollary \ref{cor:matrix-rank-nullity}(i), so $\ker(X)=\{\mathbf{0}\}$ by Proposition \ref{prop:dimension-bounds-and-basis-criteria}(ii).\\
Hence $\ker(XA)=\{\mathbf{z}\in\mathbb{F}^n:X(A\mathbf{z})=\mathbf{0}\}=\{\mathbf{z}\in\mathbb{F}^n:A\mathbf{z}=\mathbf{0}\}=\ker(A)$, and $\operatorname{rank}(XA)=\operatorname{rank}(A)$ by Corollary \ref{cor:matrix-rank-nullity}(i).\\
Since $\operatorname{rank}(Y^t)=\operatorname{rank}(Y)=n$ by Theorem \ref{thm:row-rank-equals-column-rank}, the equality $\operatorname{rank}(Y^tA^t)=\operatorname{rank}(A^t)$ holds by the preceding kernel argument with $(X,A)$ replaced by $(Y^t,A^t)$. The identity $(AY)^t=Y^tA^t$ holds by Proposition \ref{prop:transpose-reverses-products}.\\
Hence $\operatorname{rank}(AY)=\operatorname{rank}(Y^tA^t)=\operatorname{rank}(A^t)=\operatorname{rank}(A)$ by Theorem \ref{thm:row-rank-equals-column-rank} and the preceding equality.\\
Finally, $\operatorname{rank}(XAY)=\operatorname{rank}(AY)=\operatorname{rank}(A)$ by the preceding kernel argument with $A$ replaced by $AY$.
\item Set $X=P$ and $Y=Q$. Then $\operatorname{rank}(P)=m$ and $\operatorname{rank}(Q)=n$ by Corollary \ref{cor:square-matrix-invertible-iff-full-rank}, so $\operatorname{rank}(PAQ)=\operatorname{rank}(A)$ by part (i).
\end{enumerate}
\hlt{Structural Commentary}: Kernel preservation handles left factors; transposition converts right factors to the same argument.\\
\hlt{Key Insight}: Full column and row rank prevent dimensional collapse.
\end{proof}
\end{theorem}
